\documentclass{beamer}

\usepackage{cuhkbeamer}

\begin{document}

% ====================== Slide 1: TITLE SLIDE (REVISED) ======================
\begin{frame}
  \title{Economic and Social Impact of LAE}
  \author{Chiu Yu KO}
  \institute{Low Altitude Economics}
  \date{Week 9}
  \titlepage
\end{frame}


% ====================== Slide 2: FORMAT & OBJECTIVES ======================
\begin{frame}[t]
\frametitle{Learning Objectives}
By the end of this class, you should be able to:
\begin{itemize}
\item \textbf{Build} a transparent market-size estimate for a defined Hong Kong LAE use case.
\item \textbf{Distinguish} potential demand, serviceable demand, capacity, revenue, and value added.
\item \textbf{Stress-test} a business case using scenarios, sensitivity analysis, and simple probability models.
\item \textbf{Identify} who receives the benefits and who bears the costs of an LAE project.
\item \textbf{Explain} how hedonic analysis and benefit-cost analysis may inform noise, access, and equity decisions.
\end{itemize}
\end{frame}

% ====================== Slide 3: ICEBREAKER ======================
\begin{frame}[t]
\frametitle{From Commercial Value to Social Value}
\textbf{Imagine that several LAE services are operating across Hong Kong.}
\begin{itemize}
\item Users may benefit from faster delivery, inspection, emergency response, or passenger travel.
\item Nearby communities may experience noise, privacy concerns, visual effects, or unequal access to benefits.
\end{itemize}
\vspace{0.3cm}
A profitable service can still create social costs, while an unprofitable service may create important public benefits.
\begin{itemize}
\item \textbf{Quick poll (2 min):} Which matters most for public policy: operator profit, total social benefit, or how benefits and costs are distributed?
\end{itemize}
\end{frame}


% ====================== Session 1 Title Card ======================
\begin{frame}[c]
  \begin{center}
    \huge\color{cuhkpurple}\textbf{Session 1}\\
    \vspace{0.5cm}
    \Large Market Sizing, NASA Barriers \& Uncertainty Modeling
  \end{center}
\end{frame}

% ====================== Slide 4: CONCEPT PRIMER 1 - MARKET SIZING ======================
\begin{frame}[t]
\frametitle{Concept Primer: Market Sizing}
Market sizing should connect the demand analysis from Week 4 with the capacity and cost analysis from later weeks.
\begin{itemize}
\item \textbf{Potential market:} All relevant customers or trips within a clearly defined geography, period, and use case.
\item \textbf{Serviceable market:} The portion compatible with the proposed route, service level, willingness to pay, regulation, and infrastructure.
\item \textbf{Obtainable demand:} The portion the operator may realistically serve after competition and capacity constraints.
\end{itemize}
\textit{TAM, SAM, and SOM are managerial labels, not a unique NASA methodology. NASA-sponsored UAM studies use detailed market models and examine multiple constraints.}
\end{frame}

% ====================== Slide 4.5: HOW TO GET TAM ======================
\begin{frame}[t]
\frametitle{Applied Tool: Defining the Potential Market}
\textbf{Start with the decision, not the largest available number.}
\begin{enumerate}
\item Define the service, origin-destination market, customer, direction, and time period.
\item Estimate relevant existing trips or missions using transport, airport, logistics, public-service, or survey data.
\item Decide whether the project may divert existing activity, generate new activity, or change trip frequency.
\item Remove trips that cannot use the service because of geography, schedule, payload, accessibility, regulation, or other constraints.
\item Document data gaps and use ranges where evidence is weak.
\end{enumerate}
\textit{Do not sum all trips to the airport unless they are genuinely relevant to the defined corridor and service.}
\end{frame}

% ====================== Slide 5: APPLIED TOOL - MARKET FUNNEL ======================
\begin{frame}[t]
\frametitle{Applied Tool: A Market Funnel with Capacity}
\begin{block}{Demand Funnel}
\[
SAM=TAM\times\text{Eligibility Share}
\]
\[
\text{Potential Operator Demand}=SAM\times\text{Expected Choice Share}
\]
\end{block}
\begin{block}{Capacity Check}
\[
\text{Served Demand}=\min\{\text{Potential Operator Demand},\text{Reliable Capacity}\}
\]
\end{block}
\begin{itemize}
\item Choice share should reflect price, generalized cost, service quality, and competition.
\item Reliable capacity includes aircraft, seats or payload, infrastructure, operating hours, weather, and approvals.
\item Avoid calling an unsupported capture percentage ``realistic'' without evidence.
\end{itemize}
\end{frame}

% ====================== Slide 6: APPLIED TOOL - DAILY REVENUE ======================
\begin{frame}[t]
\frametitle{Applied Tool: Revenue and Contribution}
\begin{block}{Revenue}
\[
\text{Daily Revenue}=\text{Served Units}\times\text{Average Revenue per Unit}
\]
\end{block}
\begin{block}{Contribution}
\[
\text{Daily Contribution}=\text{Daily Revenue}-\text{Variable Cost of Served Units}
\]
\end{block}
\begin{itemize}
\item Use passenger-trips, deliveries, inspections, or missions consistently.
\item Revenue is not cash flow or profit and does not by itself justify infrastructure investment.
\item Compare contribution and longer-term cash flows with fixed operating cost and investment cost.
\end{itemize}
\end{frame}

% ====================== Slide 7: BACK-OF-THE-ENVELOPE - MARKET FUNNEL ======================
\begin{frame}[t,shrink=10]
\frametitle{Back-of-the-Envelope: Filtering and Capacity}
\textbf{Illustrative airport-corridor assumptions:}
\begin{columns}[T]
\begin{column}{0.48\textwidth}\begin{block}{Demand Funnel}
\begin{itemize}
\item Relevant trips: 20,000/day
\item Eligible time-sensitive segment: 10\%
\item $SAM=20{,}000\times0.10=2{,}000$
\item Expected choice share: 15\%
\item Potential operator demand: \textbf{300 passenger-trips/day}
\end{itemize}\end{block}\end{column}
\begin{column}{0.48\textwidth}\begin{block}{Capacity Check}
\begin{itemize}
\item 80 flights/day
\item Four seats per flight
\item Target load factor: 80\%
\item Reliable capacity: $80\times4\times0.80=256$
\item \textbf{Served demand: 256 passenger-trips/day}
\end{itemize}\end{block}\end{column}
\end{columns}
\textit{Takeaway: Obtainable demand is constrained by both customer choice and reliable operating capacity.}
\vfill{\tiny Illustrative assumptions; not current airport-corridor data.}
\end{frame}

% ====================== Slide 8: BACK-OF-THE-ENVELOPE - DAILY REVENUE ======================
\begin{frame}[t,shrink=10]
\frametitle{Back-of-the-Envelope: From Served Demand to Contribution}
\textbf{Assumptions:} 256 passenger-trips/day, average fare HK\$400, variable cost HK\$120/passenger-trip, and 330 operating days/year.
\begin{columns}[T]
\begin{column}{0.48\textwidth}\begin{block}{Daily Result}
\begin{itemize}
\item Revenue: $256\times400=\text{HK\$102,400}$
\item Variable cost: $256\times120=\text{HK\$30,720}$
\item \textbf{Daily contribution: HK\$71,680}
\end{itemize}\end{block}\end{column}
\begin{column}{0.48\textwidth}\begin{block}{Annual Result}
\begin{itemize}
\item Annual revenue: HK\$33.8M
\item Annual variable cost: HK\$10.1M
\item \textbf{Annual contribution: about HK\$23.7M}
\item Fixed cost and investment remain to be covered.
\end{itemize}\end{block}\end{column}
\end{columns}
\vfill{\tiny Illustrative assumptions; excludes fixed cost, financing, tax, and demand variation.}
\end{frame}

% ====================== Slide 9: CONCEPT PRIMER 2 - NASA BARRIERS ======================
\begin{frame}[t]
\frametitle{Concept Primer: Constraints Identified in UAM Market Studies}
The NASA-sponsored UAM market study examined interacting constraints rather than a universal list of five revenue haircuts.
\begin{itemize}
\item legal and regulatory requirements;
\item certification and safety;
\item infrastructure and airspace capacity;
\item weather and technology limitations;
\item public perception and community acceptance; and
\item high near-term service cost and other economic constraints.
\end{itemize}
\textbf{Managerial lesson:} Each constraint may affect demand, capacity, price, cost, timing, or probability of launch through a different channel.
\end{frame}

% ====================== Slide 10: APPLIED TOOL - NASA BARRIER HAIRCUT ======================
\begin{frame}[t]
\frametitle{Applied Tool: Constraint Mapping}
\begin{center}\small
\begin{tabular}{l|l}
\textbf{Constraint} & \textbf{Primary model channel}\\
\hline
Regulatory delay & Launch date, cost, or probability\\
Infrastructure scarcity & Reliable capacity and access\\
Public acceptance & Choice share, hours, or route\\
Weather & Serviceable capacity and reliability\\
Technology maturity & Cost, payload, range, and downtime\\
Competition & Price and operator choice share\\
\end{tabular}
\end{center}
\vspace{0.3cm}
\textit{Avoid multiplying arbitrary independent ``haircuts'' when the same effect is already reflected elsewhere. For example, weather should not reduce both operating days and revenue again unless the channels are distinct.}
\end{frame}

% ====================== Slide 11: BACK-OF-THE-ENVELOPE - BARRIER HAIRCUT ======================
\begin{frame}[t,shrink=12]
\frametitle{Back-of-the-Envelope: Explicit Constraint Channels}
\textbf{Base assumptions:} 300 potential passenger-trips/day, fare HK\$400, and 365 calendar days.
\begin{columns}[T]
\begin{column}{0.48\textwidth}\begin{block}{Demand and Capacity}
\begin{itemize}
\item Reliable capacity: 256 trips/day
\item Acceptance scenario: 90\% of potential demand
\item Adjusted demand: 270 trips/day
\item Served trips: $\min\{270,256\}=256$
\end{itemize}\end{block}\end{column}
\begin{column}{0.48\textwidth}\begin{block}{Operating Days and Revenue}
\begin{itemize}
\item Serviceable days: 330/year
\item Annual revenue: $256\times400\times330$
\item \textbf{HK\$33.8M}
\item Technology and compliance costs enter the cost model rather than another revenue haircut.
\end{itemize}\end{block}\end{column}
\end{columns}
\textit{Takeaway: Model each constraint through its actual channel and check for double counting.}
\vfill{\tiny Illustrative assumptions for classroom analysis.}
\end{frame}


% ====================== Slide 10: SESSION 2 TITLE ======================
\begin{frame}[c]
  \begin{center}
    \huge\color{cuhkpurple}\textbf{Session 2}\\
    \vspace{0.5cm}
    \Large Uncertainty, Monte Carlo Sensitivity \& Macro Impact
  \end{center}
\end{frame}

% ====================== Slide 10.5: CONCEPT PRIMER - UNCERTAINTY ======================
\begin{frame}[t]
\frametitle{Concept Primer: Uncertainty}
Uncertainty means that important inputs and outcomes cannot be known precisely when a decision is made.
\begin{itemize}
\item demand and willingness to pay;
\item approval timing and operating conditions;
\item weather and serviceable capacity;
\item technology performance and maintenance;
\item public acceptance and competition; and
\item prices, costs, and financing conditions.
\end{itemize}
\textbf{Managerial lesson:} Use a base case, downside and upside scenarios, sensitivity analysis, and probability ranges where justified. Not every uncertainty is measurable or hedgeable.
\end{frame}

% ====================== Slide 11: CONCEPT PRIMER 3 - MONTE CARLO ======================
\begin{frame}[t]
\frametitle{Concept Primer: Monte Carlo Simulation}
A Monte Carlo simulation repeatedly samples uncertain inputs from specified probability distributions and recalculates the model.
\begin{itemize}
\item It produces a distribution of possible outcomes rather than one forecast.
\item Results depend on the chosen distributions, correlations, model structure, and data quality.
\item Correlated inputs matter. Weather may affect capacity, demand, price, and cost simultaneously.
\item Useful outputs include median NPV, downside percentiles, probability of a cash shortfall, and major risk drivers.
\end{itemize}
\textit{Simulation does not make weak assumptions accurate; it makes their implications visible.}
\end{frame}

% ====================== Slide 12: APPLIED TOOL - EXPECTED VALUE ======================
\begin{frame}[t]
\frametitle{Applied Tool: Expected Value and Downside Risk}
\begin{block}{Expected Profit}
\[
E[\Pi]=\sum_s \Pr(s)\Pi_s
\]
\end{block}
\begin{itemize}
\item Expected value is the probability-weighted average across scenarios.
\item It does not reveal how severe the downside can be or whether the business has enough liquidity to survive it.
\item Also examine the probability of loss, worst credible cash shortfall, and a downside percentile such as the 10th percentile.
\item Required cash reserves depend on the timing of cash flows, financing access, risk tolerance, and recovery plan, not expected value alone.
\end{itemize}
\end{frame}

% ====================== Slide 13: BACK-OF-THE-ENVELOPE - MONTE CARLO ======================
\begin{frame}[t,shrink=11]
\frametitle{Back-of-the-Envelope: Expected Profit and Probability of Loss}
\textbf{Assumption:} Annual cost = HK\$5M.
\begin{columns}[T]
\begin{column}{0.48\textwidth}\begin{block}{Three Scenarios}
\begin{itemize}
\item Strong: 20\%, revenue HK\$10M, profit HK\$5M
\item Base: 50\%, revenue HK\$8M, profit HK\$3M
\item Weak: 30\%, revenue HK\$4M, loss HK\$1M
\end{itemize}\end{block}\end{column}
\begin{column}{0.48\textwidth}\begin{block}{Risk Summary}
\begin{itemize}
\item Expected revenue: HK\$7.2M
\item \textbf{Expected profit: HK\$2.2M}
\item \textbf{Probability of loss: 30\%}
\item Minimum cash needed in the weak scenario depends on when costs and revenue occur.
\end{itemize}\end{block}\end{column}
\end{columns}
\textit{Takeaway: A positive expected profit can coexist with a material probability of loss.}
\vfill{\tiny Illustrative scenario probabilities and values.}
\end{frame}

% ====================== Slide 14: MACRO IMPACT - HK\$350B PRIZE \& SOCIAL EQUITY ======================
\begin{frame}[t]
\frametitle{A HK\$350 Billion Projection and Distributional Questions}
A 2025 report by Our Hong Kong Foundation projects that Hong Kong's LAE-related industries could generate approximately HK\$350 billion in annual industry value added by 2050, with supporting services contributing the majority.
\begin{itemize}
\item This is a scenario-based projection, not government revenue, operator revenue, GDP automatically created, or a guaranteed outcome.
\item Value added measures the contribution of labour and capital after subtracting intermediate inputs; it should not be added directly to market revenue estimates.
\item Possible benefits include productivity, logistics, public-service performance, employment, research, finance, insurance, legal, and professional services.
\item Distribution matters: who gains access, who receives employment and business opportunities, and who bears noise, privacy, safety, land-use, or fiscal costs?
\end{itemize}
\end{frame}

% ====================== Slide 15: SESSION 3 TITLE ======================
\begin{frame}[c]
  \begin{center}
    \huge\color{cuhkpurple}\textbf{Session 3}\\
    \vspace{0.5cm}
    \Large Stakeholder Effects, Hedonic Valuation \& Equity
  \end{center}
\end{frame}

% ====================== Slide 16: CONCEPT PRIMER 4 ======================
\begin{frame}[t]
\frametitle{Concept Primer: Hedonic Valuation}
Hedonic analysis estimates how observable attributes are associated with market prices while controlling for other factors.
\begin{itemize}
\item For housing, attributes may include size, age, location, transport access, schools, amenities, and environmental conditions such as noise.
\item A credible study needs many comparable transactions and a strategy for separating noise effects from other site characteristics.
\item A simple difference between two homes does not establish the causal value of noise.
\item Property-price effects capture capitalization in housing markets, not every health, annoyance, productivity, or distributional effect.
\end{itemize}
\end{frame}

% ====================== Slide 17: STANDALONE FORMULA ======================
\begin{frame}[t]
\frametitle{Applied Tool: A Hedonic Price Model}
\begin{block}{Illustrative Specification}
\[
\ln(P_i)=\alpha+\beta N_i+\gamma A_i+\delta X_i+\varepsilon_i
\]
\end{block}
\begin{itemize}
\item $P_i$ = property price or rent.
\item $N_i$ = measured noise exposure.
\item $A_i$ = access to the relevant service or transport network.
\item $X_i$ = other property and neighbourhood characteristics.
\item $\beta$ estimates the conditional association between noise and price; causal interpretation requires stronger research design.
\end{itemize}
\textbf{Managerial use:} Inform benefit-cost analysis and mitigation, not decide community acceptance by property value alone.
\end{frame}

% ====================== Slide 18: BACK-OF-ENVELOPE (HEDONIC) ======================
\begin{frame}[t,shrink=12]
\frametitle{Back-of-the-Envelope: Illustrative Property-Value Scenarios}
\textbf{Assumption:} Baseline apartment value = HK\$8M. The percentages below are hypothetical, not estimates for drone noise in Hong Kong.
\begin{columns}[T]
\begin{column}{0.48\textwidth}\begin{block}{Higher Noise, Limited Local Access}
\begin{itemize}
\item Assumed noise effect: $-$4\%
\item Assumed access effect: 0\%
\item New value: $8M\times0.96$
\item \textbf{HK\$7.68M}
\end{itemize}\end{block}\end{column}
\begin{column}{0.48\textwidth}\begin{block}{Lower Noise, Useful Local Access}
\begin{itemize}
\item Assumed noise effect: $-$1\%
\item Assumed access effect: +3\%
\item New value: $8M\times1.02$
\item \textbf{HK\$8.16M}
\end{itemize}\end{block}\end{column}
\end{columns}
\textit{Takeaway: The arithmetic illustrates competing effects. Actual coefficients require local transaction, noise, access, and neighbourhood data.}
\end{frame}

% ====================== Slide 19: SOCIAL EQUITY & ESG ======================
\begin{frame}[t]
\frametitle{Social Equity and Distributional Analysis}
A social-impact assessment should ask who receives each benefit and who bears each cost.
\begin{itemize}
\item \textbf{Access}: affordability, geographic coverage, accessibility, digital access, and service availability.
\item \textbf{Exposure}: noise, privacy, visual impact, safety risk, construction, and land-use effects.
\item \textbf{Opportunity}: jobs, training, procurement, entrepreneurship, and supporting-service development.
\item \textbf{Public finance}: subsidies, infrastructure spending, availability payments, and alternative uses of public funds.
\item \textbf{Public services}: emergency response and medical logistics may create broad benefits, but should be evaluated against alternative delivery methods.
\end{itemize}
\textit{ESG reporting can support transparency, but a project should define measurable outcomes rather than assume that an EMS label guarantees equity.}
\end{frame}

% ====================== Slide 20: GLOBAL COMPARATORS ======================
\begin{frame}[t]
\frametitle{Comparing Social Impact Across Places}
Social impacts depend on local context and should not be ranked through broad claims.
\begin{itemize}
\item Hong Kong's dense urban environment makes access, noise exposure, land use, existing public transport, and distributional effects especially important.
\item Different GBA cities may have different trip patterns, industrial structures, land availability, public-service needs, and regulatory settings.
\item International evidence on aviation noise and property values can inform hypotheses, but coefficients should not be transferred mechanically to drone or eVTOL operations in Hong Kong.
\item A common evaluation framework can compare safety, time savings, emissions, noise, accessibility, employment, fiscal cost, and distribution across stakeholder groups.
\end{itemize}
\end{frame}

% ====================== Slide 21: EXERCISE INSTRUCTIONS ======================
\begin{frame}[t]
\frametitle{In-Class Exercise: Market, Risk, and Social Impact}
\textbf{Activity: Stress-Testing the Business and Public Case}
\begin{itemize}
\item \textbf{Format:} Groups of 4--5 with the spreadsheet template
\item \textbf{Tasks:}
\begin{enumerate}
\item define relevant trips, eligibility, expected choice share, and reliable capacity;
\item calculate served demand, revenue, variable cost, and contribution;
\item vary at least three uncertain inputs and report expected profit and probability of loss;
\item identify one benefit and one cost for each major stakeholder group; and
\item test one mitigation and state who pays for it.
\end{enumerate}
\item \textbf{Timing:} 15 minutes plus 5-minute discussion.
\end{itemize}
\end{frame}

% ====================== Slide 22: EXERCISE TEMPLATE ======================
\begin{frame}[t]
\frametitle{Exercise Template: Integrated Dashboard}
\begin{center}\small
\begin{tabular}{l|l}
\textbf{Section} & \textbf{Inputs and outputs}\\
\hline
Market & Relevant trips, eligibility, choice share\\
Capacity & Flights, capacity per flight, serviceable days\\
Economics & Price, variable cost, fixed cost, contribution\\
Uncertainty & Ranges, distributions, correlations, probability of loss\\
Social impact & Benefits, costs, affected groups, mitigation\\
\end{tabular}
\end{center}
\vspace{0.3cm}
\textit{Do not use unexplained barrier toggles. Each scenario assumption should change a specific model input and include a source or justification.}
\end{frame}

% ====================== Slide 23: KEY TAKEAWAYS ======================
\begin{frame}[t]
\frametitle{Key Takeaways}
\begin{itemize}
\item Market size must be defined by service, corridor, customer, period, choice behaviour, and reliable capacity.
\item Revenue, contribution, profit, value added, and social benefit are different concepts.
\item Constraints should enter through explicit demand, capacity, price, cost, timing, or probability channels without double counting.
\item Expected value is useful, but decision-makers also need probability of loss, downside cash needs, and key risk drivers.
\item Hedonic analysis requires comparable data and careful identification; simple assumed property discounts are scenarios, not evidence.
\item Social evaluation must consider total benefits and costs together with their distribution across users, communities, workers, taxpayers, and businesses.
\end{itemize}
\end{frame}

% ====================== Slide 24: ASSIGNMENT ======================
\begin{frame}[t]
\frametitle{Week 9 Assignment Briefing}
\textbf{Market, Risk, and Social-Impact Memo with Model (maximum 1,200 words plus spreadsheet)}
\begin{itemize}
\item \textbf{Due:} Beginning of Week 10 (group assignment)
\item \textbf{Requirements:}
\begin{enumerate}
\item define the relevant market using sourced data and transparent exclusions;
\item estimate serviceable and served demand after a reliable-capacity check;
\item calculate revenue, variable cost, contribution, and relevant fixed cost;
\item model at least three uncertain inputs, including correlations where material;
\item report expected profit, probability of loss, and one downside cash-flow measure;
\item identify benefits, costs, and distributional effects for at least four stakeholder groups; and
\item evaluate one mitigation, including its cost, beneficiary, and payer.
\end{enumerate}
\item \textbf{Grading focus:} Model consistency (30\%), evidence and assumptions (25\%), risk analysis (20\%), social-impact analysis (15\%), professionalism (10\%).
\end{itemize}
\end{frame}

% ====================== Slide 25: PREVIEW WEEK 10 ======================
\begin{frame}[t]
  \frametitle{Preview of Week 10}
  
  \begin{center}
      \Large \textbf{Sustainability Challenges in Urban Air Operations}
  \end{center}
  
  \vspace{1cm}
  \textbf{Preview question for next week:} \\
  \textit{We just quantified the impact of noise and emissions... but what is the actual mathematical cost of solving them?}
\end{frame}

% ====================== Slide 26: THANK YOU ======================
\begin{frame}[t]
  \frametitle{Thank You + Q\&A}
  
  \begin{center}
      \textbf{Next Week Preview:} \\
      \Large Green Skies: Carbon, Noise, and the Cost of Clean Tech
  \end{center}
  
  \vspace{0.5cm}
  \textbf{Pre-Class Readings Reminder (for Week 10):}
  \begin{itemize}
      \item \textit{NASA UAM Market Study (Sustainability metrics)}
      \item \textit{Our Hong Kong Foundation (ESG mandates for LAE)}
  \end{itemize}
  
  \vspace{0.5cm}
  \begin{center}
      \textbf{Questions?} \\
      \small (Final 15 minutes reserved for extended Q\&A)
  \end{center}
\end{frame}

\end{document}